\documentclass[11pt,a4paper]{article}
\usepackage[utf8]{inputenc}
\usepackage[T1]{fontenc}
\usepackage{lmodern}
\usepackage[french]{babel}
\usepackage{geometry}
\usepackage{hyperref}
\usepackage{enumitem}
\usepackage{listings}
\usepackage{longtable}
\usepackage{array}
\usepackage{seqsplit}
\geometry{margin=2.2cm}
\hypersetup{colorlinks=true,linkcolor=blue,urlcolor=blue}

\lstdefinestyle{cmd}{
  basicstyle=\ttfamily\small,
  breaklines=true,
  frame=single,
  columns=fullflexible,
  keepspaces=true
}

\title{Guide operatoire\\Parametres du workflow \emph{Manual Auto Update KB}\\et demarrage rapide du runner self-hosted}
\author{Projet STM32Cube Preprocessing}
\date{\today}

\begin{document}
\maketitle
\tableofcontents
\newpage

\section{Objet du document}
Ce document a deux objectifs:
\begin{itemize}[leftmargin=*]
  \item expliquer, champ par champ, les parametres du workflow GitHub Actions
  \texttt{Manual Auto Update KB} (\texttt{.github/workflows/manual\_auto\_update\_kb.yml}),
  declenchable manuellement depuis l'onglet Actions;
  \item fournir un guide rapide pour installer et configurer une machine de test
  equipee d'un runner self-hosted, prete a executer ce workflow.
\end{itemize}

\section{Parametres du workflow \emph{Manual Auto Update KB}}
Le workflow est declenche via \texttt{workflow\_dispatch} et propose les
parametres suivants dans l'interface GitHub Actions.

\renewcommand{\arraystretch}{1.35}
\begin{longtable}{|p{4.4cm}|p{2.2cm}|p{8.2cm}|}
\hline
\textbf{Champ} & \textbf{Type / Defaut} & \textbf{Explication} \\
\hline
\endhead
\footnotesize\texttt{\seqsplit{series}} & \footnotesize texte, defaut vide &
\footnotesize Codes de series STM32 separes par des virgules (ex: \texttt{H7,U5,WB}). Les valeurs
autorisees sont derivees de \texttt{\seqsplit{shared/config/config\_all\_series.json}} (prefixe
\texttt{STM32Cube} retire, ex: \texttt{STM32CubeH7} devient \texttt{H7}). Une valeur
vide traite toutes les series. \\
\hline
\footnotesize\texttt{mode} & \footnotesize choix, defaut \texttt{Full} &
\footnotesize Mode d'execution du script par serie. \texttt{Full} = pipeline complet
(ingestion jusqu'a la livraison et l'upload). \texttt{Prepare} = genere uniquement
les donnees ST-ready sans upload. \texttt{Upload} = upload seul, en supposant que
les fichiers sont deja prepares. \\
\hline
\footnotesize\texttt{\seqsplit{existing\_datasource\_mode}} & \footnotesize choix, defaut \texttt{Replace} &
\footnotesize Comportement si le datasource existe deja cote Alfred/KB. \texttt{Replace} = supprime
et recree (rafraichissement complet). \texttt{Add} = ajoute en complement sans
toucher a l'existant. \\
\hline
\footnotesize\texttt{\seqsplit{skip\_drivers}} & \footnotesize booleen, defaut \texttt{false} &
\footnotesize Si vrai, saute l'upload des drivers et sous-depots, pour ne traiter que les
issues/fichiers/resolveurs. \\
\hline
\footnotesize\texttt{\seqsplit{skip\_schema\_validation}} & \footnotesize booleen, defaut \texttt{false} &
\footnotesize Si vrai, desactive le garde-fou \texttt{validate\_schemas.py} avant l'upload.
Reserve au depannage, jamais en production. \\
\hline
\footnotesize\texttt{\seqsplit{continue\_on\_workflow\_error}} & \footnotesize booleen, defaut \texttt{false} &
\footnotesize Si vrai, lorsque \texttt{run\_full\_workflow} traite plusieurs depots et qu'un
depot echoue, l'execution continue avec les depots suivants au lieu de tout
arreter. \\
\hline
\footnotesize\texttt{\seqsplit{placeholder\_policy}} & \footnotesize choix, defaut \texttt{Fail} &
\footnotesize Politique appliquee si des placeholders Alfred (images/figures non resolues)
subsistent dans les donnees. \texttt{Fail} bloque le pipeline, \texttt{Warn} journalise
un avertissement et continue, \texttt{Off} ignore completement le controle. \\
\hline
\footnotesize\texttt{\seqsplit{enable\_issues\_by\_component\_flow}} & \footnotesize booleen, defaut \texttt{false} &
\footnotesize Active la generation de datasources \og issues \fg{} separes par composant, en plus
ou a la place du datasource monolithique. \\
\hline
\footnotesize\texttt{\seqsplit{component\_issues\_only}} & \footnotesize booleen, defaut \texttt{false} &
\footnotesize Si vrai, ne conserve que les datasources issues par composant et saute la version
monolithique complete. \\
\hline
\footnotesize\texttt{\seqsplit{max\_parallel}} & \footnotesize texte (nombre), defaut \texttt{'1'} &
\footnotesize Nombre maximal de jobs de series executes en parallele dans la matrice GitHub
Actions. \\
\hline
\footnotesize\texttt{\seqsplit{upload\_execution\_target}} & \footnotesize choix, defaut \texttt{self-hosted} &
\footnotesize Runner utilise pour l'etape d'upload. \texttt{self-hosted} = runner local (acces
reseau ST normal). \texttt{github-hosted} = runner cloud GitHub, reserve aux tests
sans acces reseau interne ST. \\
\hline
\end{longtable}

\section{Notions cles utilisees dans ce document}
Cette section definit les termes techniques employes ci-dessus, afin d'eviter
toute ambiguite de lecture.

\begin{itemize}[leftmargin=*]
  \item \textbf{Datasource (KB)}: unite de contenu enregistree dans la Knowledge
  Base Alfred (par exemple, les issues d'une serie, ou les fichiers d'un
  composant). Chaque datasource possede un identifiant numerique unique
  (\texttt{IssuesDatasourceId}, etc.) trace dans
  \texttt{shared/config/config\_all\_series.json}.
  \item \textbf{Alfred}: le service d'IA interne ST utilise pour enrichir le
  contenu (par exemple decrire des figures/PDF) et pour heberger la Knowledge
  Base consultee par le chatbot final.
  \item \textbf{Placeholder Alfred}: marqueur temporaire insere dans un document
  en attendant qu'Alfred fournisse une description (ex: pour une image ou un
  schema extrait d'un PDF). Un placeholder \og non resolu \fg{} signifie que la
  description finale n'a pas encore ete generee ou integree.
  \item \textbf{Runner GitHub Actions}: machine qui execute les jobs d'un
  workflow. Un runner \textbf{self-hosted} est une machine geree par l'equipe
  projet (acces reseau interne ST), alors qu'un runner \textbf{github-hosted}
  est une machine ephemere fournie par GitHub dans le cloud, sans acces au
  reseau interne ST.
  \item \textbf{Job / matrice de jobs}: un \emph{job} est une unite d'execution
  d'un workflow (ex: traiter une serie). Une \emph{matrice} lance plusieurs jobs
  en parallele, un par valeur d'une liste (ex: une serie par job). Le parametre
  \texttt{max\_parallel} limite combien de ces jobs tournent simultanement.
  \item \textbf{Schema de validation}: fichier de regles (via
  \texttt{validate\_schemas.py}) qui verifie que chaque JSON de livraison
  ST-ready respecte la structure attendue (champs obligatoires, types) avant
  d'etre uploade vers la Knowledge Base.
\end{itemize}

\section{Guide rapide: installation d'un runner self-hosted}
Objectif: installer et configurer une machine de test en dix minutes.

\subsection{Prerequis machine}
\begin{itemize}[leftmargin=*]
  \item Windows 10/11.
  \item Git installe.
  \item Python 3.12 ou superieur installe.
  \item Machine sans veille (sleep/hibernate desactive).
  \item Connexion Internet stable.
\end{itemize}

\subsection{Creation du runner GitHub}
\begin{enumerate}[leftmargin=*]
  \item Ouvrir le depot GitHub concerne.
  \item Aller dans \emph{Settings} \textrightarrow{} \emph{Actions} \textrightarrow{} \emph{Runners} \textrightarrow{} \emph{New self-hosted runner}.
  \item Choisir \emph{Windows x64}.
  \item Executer les commandes proposees par GitHub sur la machine cible.
  \item Installer le runner en tant que service et le demarrer.
\end{enumerate}

\subsection{Secrets obligatoires du depot}
Aller dans \emph{Settings} \textrightarrow{} \emph{Secrets and variables} \textrightarrow{} \emph{Actions} (onglet \emph{Secrets}, pas \emph{Variables}), puis ajouter:
\begin{itemize}[leftmargin=*]
  \item \texttt{RUNNER\_CHECK\_TOKEN} -- verifie la disponibilite du runner self-hosted (job \texttt{runner-check}).
  \item \texttt{ST\_GITHUB\_ANALYZER\_API\_KEY} -- cle prioritaire utilisee pour l'upload KB.
  \item \texttt{ST\_CHATGPT\_API\_KEY} -- cle ST ChatGPT, utilisee par le smoke test et en repli pour l'upload.
  \item \texttt{ST\_REMOTE\_USER} -- utilisateur distant pour le precheck d'authentification upload (self-hosted).
\end{itemize}
Seuls des \textbf{secrets} sont necessaires ici; aucune \emph{variable} de depot n'est utilisee par ces workflows.

\subsection{Preparation de l'environnement Python}
Depuis la racine du depot:
\begin{lstlisting}[style=cmd]
python -m venv .venv
.\.venv\Scripts\Activate.ps1
python -m pip install --upgrade pip
pip install -r requirements.txt
\end{lstlisting}

\subsection{Configuration Git sous Windows}
\begin{lstlisting}[style=cmd]
git config --system core.longpaths true
\end{lstlisting}

\subsection{Lancement du test superviseur}
\begin{enumerate}[leftmargin=*]
  \item Ouvrir l'onglet GitHub Actions.
  \item Lancer le workflow \emph{Manual Auto Update KB}.
  \item Choisir le mode \texttt{Full}.
  \item Choisir une seule serie pour le test.
\end{enumerate}

\subsection{Resultat attendu}
Ordre attendu des jobs:
\begin{enumerate}[leftmargin=*]
  \item Prepare matrix inputs
  \item Check self-hosted runner availability
  \item Preprocess <Serie>
  \item Upload <Serie>
\end{enumerate}
Controles a verifier:
\begin{itemize}[leftmargin=*]
  \item \texttt{runner\_ready = true}
  \item Alfred API key preflight = OK
  \item Upload termine sans erreur
\end{itemize}

\subsection{En cas d'echec}
\begin{itemize}[leftmargin=*]
  \item \textbf{Erreur 403 sur le controle runner}: verifier \texttt{RUNNER\_CHECK\_TOKEN}.
  \item \textbf{Cle API Alfred manquante}: ajouter \texttt{ST\_CHATGPT\_API\_KEY}
  (ou \texttt{ST\_AI\_BRIDGE\_API\_KEY} / \texttt{ST\_API\_KEY}).
  \item \textbf{Perte de communication avec le runner}: redemarrer le service runner,
  verifier le reseau ainsi que la RAM/CPU disponibles, et consulter les logs dans
  \texttt{actions-runner/\_diag}.
\end{itemize}

\section{Document complementaire}
Ce guide est concu pour etre partage avec le guide
\emph{Rapport d'equipe -- Ajout d'une nouvelle serie STM32Cube dans le pipeline}
(\texttt{new\_series\_onboarding\_and\_runner\_setup.pdf}), qui couvre l'integration
complete d'une nouvelle serie dans le pipeline et dans les workflows CI/CD.

\end{document}
